\section{Related Work}

Agent defenses reason about several interfaces around a tool call: why the
planner selected it, where its inputs originated, what the implementation may
do, and whether the represented action is authorized. We focus on the missing
validation step between execution semantics and authorization: whether the
object exposed to policy preserves the implementation-level distinctions that
policy needs.

\paragraph{Action selection and causal attribution.}
AgentDojo, ToolEmu, and AgentLAB expose indirect prompt injection in tool-rich
and long-horizon tasks~\cite{agentdojo,toolemu,agentlab}. Prompt defenses and
step classifiers assess whether a proposed action appears safe
~\cite{toolsafe,safiron}; Spotlighting delimits tool output and PromptArmor
sanitizes untrusted content~\cite{spotlighting2024,promptarmor2025}; SPIN
reframes prompt-injection detection as self-supervised classification~\cite{zhou2024spin}. CaMeL and IPIGuard preserve trusted control or track
data dependencies~\cite{camel,ipiguard}; AttriGuard and CausalArmor intervene
on observations and rerun the agent to estimate which input caused a proposed
call~\cite{attriguard,causalarmor}. These systems attribute an action to its
upstream causes. Our interventions operate at the next boundary: they vary the
call or pre-state, execute the tool, and test whether the authorization
interface preserves the resulting policy distinction.

\paragraph{Authority, provenance, and mediation.}
Agent policy systems restrict tools, actions, parameters, or capabilities
~\cite{progent,miniscope,clawguard,scopegate,agentguard}; agent-specific
mandatory access control and access-control interfaces narrow permissions at
the system and tool layers~\cite{ji2026taming,sharma2026ac4a}. A recent SoK
traces many failures to trust-authorization mismatch~\cite{shi2025soktrustauth},
while CAGE certifies authorization under typed-return uncertainty for
tool-using agents~\cite{delattre2026cage}. PACT assigns semantic
roles to arguments, while AuthGraph relates clean authority to argument
provenance~\cite{pact,authgraph}. SecureClaw and commit-time authorization place
fresh checks at the effect sink~\cite{secureclaw,committimeauth}. These systems
supply authority or enforcement mechanisms. Their precision depends on the
object presented for authorization: provenance identifies origin, call
arguments encode requests, and policy decides represented effects, but none of
these components alone establishes that the representation matches what the
implementation commits. Our enforcement model builds on complete mediation,
least privilege, confused-deputy, information-flow, and provenance principles
~\cite{saltzer1975protection,dennis1966programming,hardy1988confused,denning1976lattice,buneman2001why,cheney2009provenance}.
Contextual security models likewise separate task, action, source, and data
boundaries~\cite{siu2026formalizing}.

\paragraph{Industrial agent harnesses.}
DeepSeek Harness routes tool calls through a plugin-based pre-execution
waterfall with allow, deny, and ask decisions, monotonic guards, one-shot
approval, and an immutable logged result~\cite{deepseekharness2026}. Recent harness
engineering turns prompts into auditable contracts~\cite{ahn2026promptscontracts},
and secure plan-then-execute implementations harden agent
architectures~\cite{delrosario2025architecting}. Its
Cordis substrate formalizes revertible effects and reactive coeffects for
dynamic composition~\cite{shi2026spatiotemporal}. These systems provide an
enforcement pipeline but do not validate whether the representation exposed
to that pipeline preserves the policy-relevant committed effects of a tool
call; our registration method targets exactly that input.

\paragraph{Effect specification.}
Type-and-effect systems make effects explicit by construction
~\cite{lucassen1988effects}, and ETAS brings effect rows and compiled monitors
to an agent language~\cite{etas2026}. Contract2Tool infers planning contracts,
ContractGuard protects supplied contracts, and ToolGuardian combines
descriptions, traces, mock execution, and source analysis to vet tools
~\cite{contract2tool,contractguard,toolguardian2026}. These techniques produce,
protect, or analyze semantic descriptions. We treat such a description as a
candidate authorization interface and ask whether its induced equivalence
classes preserve the distinctions required by policy.

\paragraph{Counterexample-guided interface validation.}
Counterexample-guided abstraction refinement splits a coarse abstraction after
concrete execution falsifies it~\cite{clarke2000cegar}. Causal abstraction asks
whether interventions at a lower level are preserved by a higher-level
model~\cite{pearl2009causality,beckers2019abstracting,geiger2022inducing}. We
apply the same executable discipline to an authorization boundary. A source
intervention establishes a semantic difference, a policy oracle establishes
its security relevance, and equality under the candidate interface completes a
failure certificate. The resulting validation record makes authorization
observability policy-relative, implementation-grounded, and testable.
